\begin{center}
\begin{tikzpicture}[x=1.05cm,y=1.05cm,font=\small,>=Stealth]
\path (-2,-2.1) rectangle (9,3.6);
\coordinate (A) at (-1,.25); \coordinate (B) at (3,2.25);
\coordinate (P) at (1,-.75); \coordinate (M) at (1,1.25);
\coordinate (N) at (1,.25); \coordinate (H) at (.2,.85);
\fill[blue!12] (P)--(A)--(M)--cycle;
\fill[orange!20] (P)--(M)--(B)--cycle;
\draw[->,gray] (-1.8,0)--(3.8,0) node[right] {$x$};
\draw[->,gray] (0,-1.35)--(0,3.2) node[above] {$y$};
\draw[blue!70!black,thick,domain=-1.6:3.45,samples=100] plot (\x,{\x*\x/4});
\draw[gray!80,domain=-1.7:1.5] plot (\x,{-.5*\x-.25});
\draw[gray!80,domain=.7:3.4] plot (\x,{1.5*\x-2.25});
\draw[very thick] (A)--(B);
\draw[thick] (A)--(P)--(B);
\draw[orange!80!black,very thick] (P)--(M);
\foreach \pt/\pos in {B/above left,M/above left}
  \fill (\pt) circle (1.6pt) node[\pos] {$\pt$};
\fill (A) circle (1.6pt); \node[anchor=east] at (-1.3,-.35) {$A$};
\fill (P) circle (1.6pt); \node[anchor=west] at (1.4,-.65) {$P$};

\fill[blue!70!black] (N) circle (1.7pt) node[right] {$N$};
\draw[thick] (.88,-.25)--(1.12,-.25) (.88,.75)--(1.12,.75);

\node[anchor=north west,align=left,text width=4.7cm] at (4.1,3.1) {\textbf{平分的是线段 $PM$}\\[8pt]
$x_N=x_P=x_M$\\[8pt]
$y_N=\dfrac{(x_1+x_2)^2}{8p}$\\[10pt]
$PN=NM=\dfrac{PM}{2}$\\[8pt]
$P$ 是切线交点；\\$N$ 才是同横坐标的曲线点。\\[8pt]
两段曲线弧不据此等长。};
\end{tikzpicture}
\end{center}
